% Part III -- Reusable and Reliable Programs
% Chapter 10: Files and Paths

\h{Files and Paths}
\label{ch:10}

Variables disappear when a program ends. Files let a later run recover saved
text or structured data. This chapter begins with small text files, then adds
paths, line-by-line processing, CSV, and JSON.

\begin{NoteBox}
\textbf{Prerequisites.} You should understand strings, collections, loops,
functions, and specific exception handling.
\end{NoteBox}

\begin{tcolorbox}[title=Learning outcomes,colframe=orange,colback=white,arc=2mm]
After completing this chapter, you should be able to:
\begin{itemize}
    \item create platform-independent paths with \c{pathlib.Path};
    \item write, read, and append UTF-8 text;
    \item use \c{with} so files close automatically;
    \item process a text file one line at a time;
    \item handle missing files with a specific exception;
    \item read and write CSV rows; and
    \item save and restore dictionaries with JSON.
\end{itemize}
\end{tcolorbox}

\hh{A path identifies a location}

A path describes where a file or folder is located. The \c{pathlib} module
builds paths without requiring you to type operating-system-specific separators.

\begin{lstlisting}[
    language=python,
    style=block,
    caption={Build and inspect a path},
    label={c10_first_path},
    captionpos=b
]
from pathlib import Path

data_folder = Path("study_data")
file_path = data_folder / "rooms.txt"

print(data_folder)
print(file_path)
print(file_path.name)
print(file_path.suffix)
\end{lstlisting}

The slash operator joins path parts. It does not open or create the file.

\bookfigure[0.82\textwidth]{Figures/Generated/ch10_path_tree.pdf}
{A relative path is resolved from the working folder by following a sequence of child folders and a final file name.}
{fig:ch10-path-tree}

Read a relative path as a traversal algorithm: begin at the working folder,
follow each component in order, and stop at the described location. A path can
be valid as a description even before that file exists.

\bookfigure[0.9\textwidth]{Figures/Generated/ch10_file_pipeline.pdf}
{A dependable file workflow separates path construction, opening, processing, and closing.}
{fig:ch10-file-pipeline}

\begin{SyntaxBox}{Path expressions build locations without touching the disk}
In \c{Path("study_data") / "rooms.txt"}, \c{Path(...)} creates a path object
and the slash joins another path component. The expression describes a
location. A later method such as \c{mkdir()}, \c{write_text()}, or \c{open()}
performs the filesystem action.
\end{SyntaxBox}

\begin{TryBox}{Inspect a path before creating anything}
Create a path for \c{study_data/measurements.csv}. Print its \c{name},
\c{suffix}, and \c{parent}. Predict the results first. Then create only the
parent folder with \c{mkdir(exist_ok=True)} and check \c{exists()} for both
the folder and the still-unwritten file.
\end{TryBox}

\hh{Create a folder when needed}

Use \c{mkdir()} before writing to a folder that may not exist.

\begin{lstlisting}[
    language=python,
    style=block,
    caption={Create a study data folder safely},
    label={c10_create_folder},
    captionpos=b
]
from pathlib import Path

data_folder = Path("study_data")
data_folder.mkdir(exist_ok=True)

print(data_folder.exists())
print(data_folder.is_dir())
\end{lstlisting}

The option \c{exist_ok=True} prevents an error when the folder already exists.

\hh{Write a small text file}

For short text, \c{Path.write_text()} is concise. Opening in write mode replaces
any existing content at that path.

\begin{lstlisting}[
    language=python,
    style=block,
    caption={Write three lines of UTF-8 text},
    label={c10_write_text},
    captionpos=b
]
from pathlib import Path

data_folder = Path("study_data")
data_folder.mkdir(exist_ok=True)
file_path = data_folder / "rooms.txt"

content = "Lobby,24.0\nOffice,18.0\nCafe,31.0\n"
file_path.write_text(content, encoding="utf-8")

print(f"Wrote: {file_path}")
\end{lstlisting}

The sequence \c{"\textbackslash n"} represents a line break.

\begin{NoteBox}
\textbf{Stop and predict.} What happens to the earlier content if
\c{write_text()} is called again with different text?
\end{NoteBox}

\hh{Read a small text file}

The matching \c{read_text()} method returns the entire file as one string.

\begin{lstlisting}[
    language=python,
    style=block,
    caption={Create and read a small text file},
    label={c10_read_text},
    captionpos=b
]
from pathlib import Path

data_folder = Path("study_data")
data_folder.mkdir(exist_ok=True)
file_path = data_folder / "rooms.txt"
file_path.write_text("Lobby,24.0\nOffice,18.0\n", encoding="utf-8")

content = file_path.read_text(encoding="utf-8")
print(content)
\end{lstlisting}

Reading the entire file is convenient for small data. Large files should be
processed one line at a time.

\hh{Use with and open}

The \c{open()} function provides more control. A \c{with} block closes the file
automatically, even if an exception occurs.

\begin{lstlisting}[
    language=python,
    style=block,
    caption={Write text with a context manager},
    label={c10_with_open_write},
    captionpos=b
]
from pathlib import Path

data_folder = Path("study_data")
data_folder.mkdir(exist_ok=True)
file_path = data_folder / "notes.txt"

with open(file_path, "w", encoding="utf-8") as file:
    file.write("First note\n")
    file.write("Second note\n")

print(file_path.read_text(encoding="utf-8"))
\end{lstlisting}

The name \c{file} is available inside the block. After the block, the file has
been closed.

\hh{Choose a file mode}

The mode states what the program intends to do.

\begin{center}
\begin{tabular}{ll}
\textbf{Mode} & \textbf{Meaning} \\
\c{"r"} & Read an existing file \\
\c{"w"} & Write and replace existing content \\
\c{"a"} & Append new content at the end \\
\c{"x"} & Create a new file and fail if it exists \\
\c{"rb"}, \c{"wb"} & Read or write binary data \\
\end{tabular}
\end{center}

This chapter focuses on text modes. Images and PDFs are binary files and need
binary mode when their bytes are processed directly.

\begin{ExplainBox}{File modes declare consequences}
Mode \c{"r"} requires existing content, \c{"a"} preserves it, and \c{"w"}
replaces it. Treat the mode as part of the program's data policy: decide
whether earlier content should be required, retained, or deliberately
overwritten before opening the file.
\end{ExplainBox}

\hh{Append without replacing}

Append mode preserves existing content and adds new text at the end.

\begin{lstlisting}[
    language=python,
    style=block,
    caption={Append one line to an existing file},
    label={c10_append_text},
    captionpos=b
]
from pathlib import Path

data_folder = Path("study_data")
data_folder.mkdir(exist_ok=True)
log_path = data_folder / "study_log.txt"
log_path.write_text("Chapter 9 complete\n", encoding="utf-8")

with open(log_path, "a", encoding="utf-8") as file:
    file.write("Chapter 10 started\n")

print(log_path.read_text(encoding="utf-8"))
\end{lstlisting}

Remember to include a line break if each entry should appear on its own line.

\hh{Process one line at a time}

Iterating over a file avoids loading all of it into memory.

\begin{lstlisting}[
    language=python,
    style=block,
    caption={Read and process one line at a time},
    label={c10_line_iteration},
    captionpos=b
]
from pathlib import Path

data_folder = Path("study_data")
data_folder.mkdir(exist_ok=True)
file_path = data_folder / "areas.txt"
file_path.write_text("24.0\n35.5\n18.0\n", encoding="utf-8")

total_m2 = 0.0
count = 0

with open(file_path, "r", encoding="utf-8") as file:
    for line in file:
        area_m2 = float(line.strip())
        total_m2 += area_m2
        count += 1

print(f"Count: {count}")
print(f"Total: {total_m2:.1f} m2")
\end{lstlisting}

The \c{strip()} method removes the line break and surrounding whitespace before
conversion.

\hh{Split simple delimited text}

The \c{split()} method is suitable for controlled data with a simple delimiter.

\begin{lstlisting}[
    language=python,
    style=block,
    caption={Split two fields on a controlled delimiter},
    label={c10_split_lines},
    captionpos=b
]
from pathlib import Path

data_folder = Path("study_data")
data_folder.mkdir(exist_ok=True)
file_path = data_folder / "room_areas.txt"
file_path.write_text("Lobby|24.0\nOffice|18.0\n", encoding="utf-8")

with open(file_path, "r", encoding="utf-8") as file:
    for line in file:
        room_name, area_text = line.strip().split("|")
        area_m2 = float(area_text)
        print(f"{room_name}: {area_m2:.1f} m2")
\end{lstlisting}

Do not parse general CSV files with \c{split(",")}. Quoted fields may contain
commas. Use the \c{csv} module instead.

\hh{Handle a missing file}

Reading a path that does not exist raises \c{FileNotFoundError}.

\begin{lstlisting}[
    language=python,
    style=block,
    caption={Handle a missing text file},
    label={c10_missing_file},
    captionpos=b
]
from pathlib import Path

file_path = Path("study_data") / "missing.txt"

try:
    content = file_path.read_text(encoding="utf-8")
except FileNotFoundError:
    print(f"Could not find: {file_path}")
else:
    print(content)
\end{lstlisting}

Catch the specific problem and include the path in the message. Other errors
remain visible for diagnosis.

\hh{Check a path before an optional action}

An existence check is appropriate when either state is a normal possibility.

\begin{lstlisting}[
    language=python,
    style=block,
    caption={Choose based on whether an optional file exists},
    label={c10_path_check},
    captionpos=b
]
from pathlib import Path

file_path = Path("study_data") / "optional_notes.txt"

if file_path.exists():
    print(file_path.read_text(encoding="utf-8"))
else:
    print("No optional notes were saved.")
\end{lstlisting}

For a required file, trying the operation and handling
\c{FileNotFoundError} often provides a clearer result.

\hh{Write CSV data}

CSV stores rows and columns as text. Use \c{csv.writer()} so commas and
quotation marks are handled correctly.

\begin{lstlisting}[
    language=python,
    style=block,
    caption={Write rows with the csv module},
    label={c10_csv_write},
    captionpos=b
]
import csv
from pathlib import Path

data_folder = Path("study_data")
data_folder.mkdir(exist_ok=True)
csv_path = data_folder / "rooms.csv"

rows = [
    ["name", "area_m2"],
    ["Lobby", 24.0],
    ["Studio, North", 35.5],
    ["Office", 18.0]
]

with open(csv_path, "w", newline="", encoding="utf-8") as file:
    writer = csv.writer(file)
    writer.writerows(rows)

print(csv_path.read_text(encoding="utf-8"))
\end{lstlisting}

The \c{newline=""} argument prevents extra blank rows on some systems.

\hh{Read CSV rows as dictionaries}

\c{csv.DictReader()} uses the header row as dictionary keys.

\begin{lstlisting}[
    language=python,
    style=block,
    caption={Read named CSV columns},
    label={c10_csv_read},
    captionpos=b
]
import csv
from pathlib import Path

data_folder = Path("study_data")
data_folder.mkdir(exist_ok=True)
csv_path = data_folder / "rooms.csv"

csv_path.write_text(
    "name,area_m2\nLobby,24.0\nOffice,18.0\n",
    encoding="utf-8"
)

total_m2 = 0.0
with open(csv_path, "r", newline="", encoding="utf-8") as file:
    reader = csv.DictReader(file)
    for row in reader:
        total_m2 += float(row["area_m2"])

print(f"Total: {total_m2:.1f} m2")
\end{lstlisting}

Values read from CSV are strings. Convert numeric fields explicitly.

\hh{Save nested data with JSON}

JSON represents dictionaries, lists, strings, numbers, booleans, and null
values in a text format.

\begin{lstlisting}[
    language=python,
    style=block,
    caption={Write a dictionary as formatted JSON},
    label={c10_json_write},
    captionpos=b
]
import json
from pathlib import Path

data_folder = Path("study_data")
data_folder.mkdir(exist_ok=True)
json_path = data_folder / "settings.json"

settings = {
    "units": "metric",
    "rounding_digits": 2,
    "visible_layers": ["walls", "doors", "rooms"]
}

with open(json_path, "w", encoding="utf-8") as file:
    json.dump(settings, file, indent=2)

print(json_path.read_text(encoding="utf-8"))
\end{lstlisting}

The \c{indent=2} option makes the file easier for people to read.

\hh{Restore JSON data}

Use \c{json.load()} to convert JSON text back into Python values.

\begin{lstlisting}[
    language=python,
    style=block,
    caption={Read JSON into a dictionary},
    label={c10_json_read},
    captionpos=b
]
import json
from pathlib import Path

data_folder = Path("study_data")
data_folder.mkdir(exist_ok=True)
json_path = data_folder / "settings.json"
json_path.write_text(
    '{"units": "metric", "rounding_digits": 2}',
    encoding="utf-8"
)

with open(json_path, "r", encoding="utf-8") as file:
    settings = json.load(file)

print(settings["units"])
print(settings["rounding_digits"])
\end{lstlisting}

Invalid JSON raises \c{json.JSONDecodeError}. Handle it when a file may be
edited manually or come from an uncertain source.

\hh{Put file operations in focused functions}

Small functions separate file access from the rest of the program.

\begin{lstlisting}[
    language=python,
    style=block,
    caption={Separate saving and loading into functions},
    label={c10_file_functions},
    captionpos=b
]
from pathlib import Path

def save_lines(file_path, lines):
    file_path.parent.mkdir(parents=True, exist_ok=True)
    text = "\n".join(lines) + "\n"
    file_path.write_text(text, encoding="utf-8")

def load_lines(file_path):
    text = file_path.read_text(encoding="utf-8")
    return text.splitlines()

path = Path("study_data") / "room_names.txt"
save_lines(path, ["Lobby", "Office", "Cafe"])
room_names = load_lines(path)

print(room_names)
\end{lstlisting}

The caller chooses the path and data. Each function performs one clear task.

\hh{Common mistakes}

\begin{itemize}
    \item \textbf{Using write mode unintentionally}: \c{"w"} replaces prior content.
    \item \textbf{Forgetting encoding}: specify \c{encoding="utf-8"} for text.
    \item \textbf{Forgetting line breaks}: \c{write()} does not add them automatically.
    \item \textbf{Building paths by joining strings}: prefer \c{Path} and the slash operator.
    \item \textbf{Reading a large file all at once}: iterate over its lines.
    \item \textbf{Parsing CSV with split()}: use the \c{csv} module.
    \item \textbf{Forgetting numeric conversion}: CSV values are read as strings.
    \item \textbf{Catching every exception}: handle only expected file problems.
\end{itemize}

\hh{Practice ladder}

\hhh{Level 0 -- Read and predict}

\begin{lstlisting}[
    language=python,
    style=block,
    caption={Predict two path properties},
    label={c10_practice_level0},
    captionpos=b
]
from pathlib import Path

path = Path("study_data") / "notes.txt"
print(path.name)
print(path.suffix)
\end{lstlisting}

\hhh{Level 1 -- Modify}

Change the folder and filename. Display the parent folder, filename, stem, and
suffix.

\hhh{Level 2 -- Complete}

Use the starter code to create a folder. Then add statements that write the
three material names on separate lines, read them back, and display the text.

\begin{lstlisting}[
    language=python,
    style=block,
    caption={Starter code for writing and reading three lines},
    label={c10_practice_level2},
    captionpos=b
]
from pathlib import Path

folder = Path("study_data")
folder.mkdir(exist_ok=True)
path = folder / "materials.txt"

# Add one write statement and one read statement here.
\end{lstlisting}

\hhh{Level 3 -- Apply}

Write a CSV file with columns \c{name} and \c{area_m2}. Add three data rows.
Read it with \c{csv.DictReader()} and display the total area.

\textbf{Hints}
\begin{itemize}
    \item \textbf{Hint 1}: Open the file with \c{newline=""}.
    \item \textbf{Hint 2}: Write the header before the data rows.
    \item \textbf{Hint 3}: Convert \c{row["area_m2"]} with \c{float()}.
\end{itemize}

\hhh{Level 4 -- Challenge}

Write \c{save_settings(path, settings)} and \c{load_settings(path)} using JSON.
Handle a missing file in the caller and display a helpful message.

\hh{Practice solutions}

\textbf{Level 0:} The outputs are \c{notes.txt} and \c{.txt}.

\textbf{Level 2}

\begin{lstlisting}[
    language=python,
    style=block,
    caption={One solution to Level 2},
    label={c10_solution_level2},
    captionpos=b
]
from pathlib import Path

folder = Path("study_data")
folder.mkdir(exist_ok=True)
path = folder / "materials.txt"

path.write_text("Wood\nSteel\nConcrete\n", encoding="utf-8")
text = path.read_text(encoding="utf-8")
print(text)
\end{lstlisting}

\textbf{Level 3}

\begin{lstlisting}[
    language=python,
    style=block,
    caption={One solution to Level 3},
    label={c10_solution_level3},
    captionpos=b
]
import csv
from pathlib import Path

folder = Path("study_data")
folder.mkdir(exist_ok=True)
path = folder / "practice_rooms.csv"

with open(path, "w", newline="", encoding="utf-8") as file:
    writer = csv.writer(file)
    writer.writerow(["name", "area_m2"])
    writer.writerow(["Lobby", 24.0])
    writer.writerow(["Office", 18.0])
    writer.writerow(["Cafe", 31.0])

total_m2 = 0.0
with open(path, "r", newline="", encoding="utf-8") as file:
    reader = csv.DictReader(file)
    for row in reader:
        total_m2 += float(row["area_m2"])

print(f"Total: {total_m2:.1f} m2")
\end{lstlisting}

\textbf{Level 4}

\begin{lstlisting}[
    language=python,
    style=block,
    caption={One solution to Level 4},
    label={c10_solution_level4},
    captionpos=b
]
import json
from pathlib import Path

def save_settings(path, settings):
    path.parent.mkdir(parents=True, exist_ok=True)
    with open(path, "w", encoding="utf-8") as file:
        json.dump(settings, file, indent=2)

def load_settings(path):
    with open(path, "r", encoding="utf-8") as file:
        return json.load(file)

path = Path("study_data") / "practice_settings.json"
settings = {"units": "metric", "digits": 2}
save_settings(path, settings)

try:
    restored = load_settings(path)
except FileNotFoundError:
    print(f"Could not find: {path}")
else:
    print(restored)
\end{lstlisting}

\hh{Chapter checkpoint}

\begin{enumerate}
    \item What does a path represent?
    \item Does creating a \c{Path} object create the file?
    \item What does write mode do to existing content?
    \item Why use a \c{with} block for an open file?
    \item Which mode appends content?
    \item Why process a large file line by line?
    \item Which exception reports a missing file?
    \item Why should general CSV data be handled by the \c{csv} module?
    \item What Python structures map naturally to JSON?
    \item Why specify UTF-8 encoding?
\end{enumerate}

\hh{Checkpoint answers}

\begin{enumerate}
    \item A file or folder location.
    \item No. It only represents the location.
    \item It replaces the existing content.
    \item The file closes automatically when the block ends.
    \item \c{"a"}.
    \item It limits memory use.
    \item \c{FileNotFoundError}.
    \item It correctly handles quoted commas and other CSV rules.
    \item Dictionaries, lists, strings, numbers, booleans, and null values.
    \item It gives text a consistent character encoding across systems.
\end{enumerate}

\begin{tcolorbox}[title=Part III summary,colframe=orange,colback=white,arc=2mm]
You can organize code with functions, select useful built-ins, handle expected
errors, copy mutable data intentionally, and store information in text, CSV,
and JSON files. These skills support the numerical and data-analysis chapters
that follow.
\end{tcolorbox}

